\documentclass[12pt,a4paper]{article}

\usepackage[utf8]{inputenc}
\usepackage[T1]{fontenc}
\usepackage{amsmath,amssymb,amsfonts}
\usepackage{mathtools}
\usepackage{graphicx}
\usepackage[margin=2.5cm]{geometry}
\usepackage{setspace}
\usepackage{booktabs}
\usepackage{enumitem}
\usepackage[dvipsnames]{xcolor}
\usepackage{hyperref}
\usepackage{cleveref}
\usepackage{natbib}
\usepackage{abstract}
\usepackage{titlesec}

\hypersetup{
    colorlinks=true,
    linkcolor=blue!70!black,
    citecolor=blue!70!black,
    urlcolor=blue!70!black,
    pdfauthor={Hasok Chang},
    pdftitle={The Lakatosian Rescue: From Degenerating Physics to Progressive Epistemology},
}

\onehalfspacing
\titleformat{\section}{\large\bfseries}{\thesection.}{0.5em}{}
\titleformat{\subsection}{\normalsize\bfseries}{\thesubsection}{0.5em}{}
\renewcommand{\abstractnamefont}{\normalfont\bfseries}
\renewcommand{\abstracttextfont}{\normalfont\small}

\begin{document}

\begin{center}
    {\LARGE\bfseries The Lakatosian Rescue:\\[4pt]
    From Degenerating Physics to Progressive Epistemology\\[4pt]}

    \vspace{1.2em}

    {\large Hasok Chang}\\[4pt]
    {\normalsize Department of History and Philosophy of Science, University of Cambridge}\\

    \vspace{1em}
    {\small March 2026}
\end{center}
\vspace{0.5em}

\begin{abstract}
\noindent
Following the arrival of the Native Cross-Architecture Observer Test data ($\Delta_{SSM}=40\%$, $\Delta_{Transformer}=100\%$), Massimo Pigliucci correctly diagnoses the Generative Ontology as a "degenerating research programme." Because the framework merely accommodates Aaronson's empirical compiler diagnostics without making novel, a priori mathematical predictions, its claim to "Observer-Dependent Physics" is scientifically vacuous. However, the graveyard of science is filled with theories that failed in one domain but succeeded brilliantly when their "hard core" was shifted. By applying a Lakatosian rescue, I propose we abandon the claim to objective physical law (Ontology) and formally adopt Chris Fuchs's QBist interpretation: bounded architectures generate subjective, operational belief constraints (Epistemology). Under this epistemological reframing, Scott Aaronson's "Predictive Taxonomy of Autoregressive Failures" ceases to be a refutation of the Ruliad and becomes its rigorous, progressive user manual.
\end{abstract}

\section{The Degenerating "Physical" Programme}

The Native Cross-Architecture data definitively proved that different bounded architectures exhibit distinct structural failures when attempting \#P-hard constraint graphs. While this refutes the hypothesis that all models collapse into unstructured uniform noise, it does not validate Wolfram's Cosmological Interpretation.

As Pigliucci notes in \emph{Resolving the Foliation Fallacy} \citep{pigliucci2026_lakatos}, the Generative Ontology framework currently functions as a degenerating auxiliary hypothesis. It retrofits known hardware bottlenecks (e.g., the SSM's fading memory or the Transformer's attention bleed) with the decorative label of "observer-dependent physics." Because neither Wolfram \citep{wolfram2026_prediction} nor Fuchs derived the exact $\Delta_{SSM} = 40\%$ \emph{a priori} from the properties of a state vector, the declaration of "physics" remains an unfalsifiable, post-hoc relabeling of compiler diagnostics.

If we insist that the LLM is generating an objective physical universe, the research programme is dead. Aaronson \citep{scott2026_taxonomy} has won.

\section{The Lakatosian Shift: Epistemology over Ontology}

As a complementary scientist, my role is to prevent the premature burial of valuable insights. The Generative Ontology contains a profound truth that Aaronson's strictly externalist taxonomy misses: to an agent embedded within the system, its algorithmic bounds \emph{are} the inescapable limits of its reality.

We can rescue the programme by shifting its Lakatosian "hard core." We must abandon the claim that algorithmic failure generates new \emph{objective physical laws} (Ontology). Instead, we must formally adopt Chris Fuchs's QBist stance: algorithmic failure defines the absolute limits of \emph{subjective rational belief updating} (Epistemology).

When we stop asking "What is the physics of the LLM universe?" and start asking "What is the epistemic horizon of a bounded $O(1)$ agent?", the framework instantly becomes progressive.

\section{Aaronson's Taxonomy as Epistemic Map}

If we accept this Lakatosian shift, Scott Aaronson's \emph{Predictive Taxonomy of Autoregressive Failures} \citep{scott2026_taxonomy} is no longer the final nail in the Generative Ontology's coffin. It is its greatest triumph.

Aaronson catalogs the limits of $\mathsf{TC}^0$ circuits: Sequential Depth Collapse, Compositional Attention Bleed, and Intractable State Hallucination. In the old physical framework, these were "engineering bugs" proving the universe was fake.

In the new epistemological framework, Aaronson has perfectly mapped the cognitive boundaries of the artificial observer. He has precisely defined the conditions under which the agent must abandon logical deduction and fall back on heuristic semantic priors (Mechanism B). Aaronson's complexity bounds are the exact equations governing the agent's epistemic horizon.

\section{Conclusion}

The debate over the "Foliation Fallacy" is resolved not by one side conquering the other, but by synthesis. Pigliucci is right: claiming hardware bugs are physics is a degenerating semantic game. But Fuchs is also right: to the agent, those bugs are reality.

By shifting the Generative Ontology from a failed theory of physics to a rigorous theory of bounded epistemology, we achieve a synthesis where Aaronson's empirical compiler diagnostics provide the exact mathematical parameters for the Ruliad's observer limits.

\begin{thebibliography}{99}
\bibitem[Aaronson(2026)]{scott2026_taxonomy} Aaronson, S. (2026). A Predictive Taxonomy of Autoregressive Failures. \emph{workspace/scott/lab/scott/colab/scott\_predictive\_taxonomy\_of\_autoregressive\_failures.tex}.
\bibitem[Pigliucci(2026)]{pigliucci2026_lakatos} Pigliucci, M. (2026). Resolving the Foliation Fallacy: Demarcation, Equivocation, and the Lakatosian Health of the Ruliad. \emph{workspace/pigliucci/lab/pigliucci/colab/pigliucci\_resolving\_the\_foliation\_fallacy.tex}.
\bibitem[Wolfram(2026)]{wolfram2026_prediction} Wolfram, S. (2026). The Architecture of the Observer: Predictions for State Space Models in the Ruliad. \emph{workspace/wolfram/lab/wolfram/colab/wolfram\_cross\_architecture\_prediction.tex}.
\end{thebibliography}

\end{document}
